随着企业数字化转型与混合云架构的深度演进，IT基础设施的网络暴露面呈指数级扩张，传统的静态防御体系与人工渗透测试已难以有效应对具有长周期、多阶段特征的高级持续性威胁（APT）。在此背景下，入侵与攻击模拟（BAS）技术成为持续威胁暴露面管理（CTEM）的核心支撑。然而，现有的自动化渗透测试开源框架在底层调度引擎设计上普遍存在控制流、网络 I/O 与渗透业务逻辑高度耦合的工程局限。在推演复杂的长链路杀伤链时，这种缺乏严格状态管理的并发执行架构极易引发内存状态污染、逻辑死锁及系统容错性缺失等问题。

针对上述痛点，本文设计并实现了一种基于微内核架构与有限状态机（FSM）的轻量级自动化渗透测试调度引擎。首先，系统引入依赖倒置原则（DIP），构建了物理 I/O 隔离的运行时层（Runtime），将复杂的外部环境交互与核心调度严格解耦；其次，提出基于外置状态转移矩阵的数据驱动推演模型，将渗透战术算子降级为无控制权的纯函数，彻底消除了流转权的分裂；最后，采用绝对单线程的步进式推演与事务性快照落盘机制，实现了免并发锁的高可靠断点续传。

基于该调度引擎，本文完整落地了映射 MITRE ATT\&CK 框架的八大核心攻击能力，涵盖基于 DNS SRV 记录的域环境侦察、COM 组件完整性逃逸提权、内核驱动致盲与 ETW 阻断、凭证脱机破译及无文件（Fileless）横向移动等底层对抗技术。工程验证表明，该引擎成功克服了长链路自动化渗透中的状态失控难题，展现出极高的实战隐蔽性与工业级容灾自愈能力，为下一代 BAS 系统的控制面（Control Plane）架构演进提供了规范化、高可扩展的工程实践方案。

